\documentclass[11pt,a4paper]{article}
\usepackage[utf8]{inputenc}
\usepackage[T1]{fontenc}
\usepackage[italian]{babel}
\usepackage{Alegreya}
\usepackage[margin=2.5cm]{geometry}
\usepackage{titlesec}
\usepackage{enumitem}
\usepackage{xcolor}

% Configurazione titoli
\titleformat{\section}{\Large\bfseries\color{black}}{\thesection}{1em}{}
\titleformat{\subsection}{\large\bfseries\color{gray}}{\thesubsection}{1em}{}
\titleformat{\subsubsection}{\normalsize\bfseries}{\thesubsubsection}{1em}{}

% Configurazione spaziatura
\setlength{\parindent}{0pt}
\setlength{\parskip}{0.8em}

% Rimuovi numerazione pagine e sezioni per il frontespizio
\pagenumbering{gobble}

\begin{document}

% FRONTESPIZIO
\begin{center}
    {\Large\textbf{CONSERVATORIO DI MUSICA BENEDETTO MARCELLO -- VENEZIA}}

    \vspace{1cm}

    {\large\textbf{DOTTORATO DI RICERCA AFAM IN\\MUSICA, PERFORMANCE E INNOVAZIONE TECNOLOGICA\\XLI CICLO}}

    \vspace{1.5cm}

    {\LARGE\textbf{VOLONTA' INTERPRETANTE}}

    \vspace{1cm}

    \textbf{Candidato:} Alice Cortegiani

    \vspace{0.5cm}

    \textbf{Parole chiave:} Fenomenologia, Ermeneutica, Prassi interpretativa, parola4, parola5

    \vspace{1cm}

    Anno Accademico 2025/2026
\end{center}

\newpage
\pagenumbering{arabic}

% CORPO DEL PROGETTO
\section{Descrizione del progetto}

\subsection{Obiettivi e focus della ricerca}
[Descrivere chiaramente l'oggetto di studio, gli obiettivi principali e come il progetto si inserisce nelle tematiche del dottorato: performance, innovazione tecnologica, heritage veneziano, neuroscienze, inclusione]

\subsection{Modalità operative e metodologie d'indagine}
\textbf{Approcci metodologici:}
\begin{itemize}[leftmargin=1cm]
    \item Specificare metodologie: artistic research, performance studies, analisi neuroscientifiche, etc.
    \item Metodologie di ricerca archivistica o di campo
\end{itemize}

\textbf{Tecnologie utilizzate:}
\begin{itemize}[leftmargin=1cm]
    \item Software specifici
    \item Hardware e strumenti di registrazione/analisi
    \item Strumentazione tecnica specialistica
\end{itemize}

\subsection{Tipologia delle fonti}
\begin{description}[leftmargin=2cm]
    \item[Fonti primarie:] [archivi, manoscritti, registrazioni originali]
    \item[Fonti secondarie:] [letteratura specialistica, studi precedenti]
    \item[Fonti multimediali:] [video, audio, software, installazioni]
\end{description}

\subsection{Necessità di viaggi e sopralluoghi}
[Specificare spostamenti necessari per accesso ad archivi, istituzioni, collaborazioni]

\subsection{Difficoltà possibili e soluzioni previste}
\textbf{Difficoltà identificate:}
\begin{itemize}[leftmargin=1cm]
    \item Ostacoli logistici
    \item Problematiche tecniche
    \item Limitazioni economiche
    \item Problemi di accesso alle fonti
\end{itemize}

\textbf{Soluzioni proposte:}
[Descrivere strategie alternative e approcci per superare gli ostacoli identificati]

\section{Stato dell'arte}

\subsection{Inserimento negli studi pregressi}
[Inquadrare il progetto nel panorama della ricerca attuale, evidenziando gli studi fondamentali di riferimento e le lacune conoscitive da colmare]

\subsection{Esperienza del candidato}
[Specificare se il progetto rappresenta un approfondimento di ricerche precedenti condotte dal candidato o l'avvio di un percorso di ricerca completamente nuovo]

\subsection{Posizionamento innovativo}
[Chiarire in che modo il progetto si distingue e si differenzia da quanto già esistente nel panorama scientifico e artistico]

\section{Risultati attesi}

\subsection{Innovazioni rispetto allo stato dell'arte attuale}
\begin{itemize}[leftmargin=1cm]
    \item \textbf{Contributi teorici:} [nuovi modelli interpretativi, framework concettuali]
    \item \textbf{Sviluppi tecnologici:} [software, hardware, metodologie innovative]
    \item \textbf{Pratiche performative:} [nuove modalità esecutive, performance ibride]
    \item \textbf{Prodotti artistici:} [composizioni, installazioni, progetti multimediali]
\end{itemize}

\subsection{Competenze e punti di forza del candidato}
\begin{description}[leftmargin=2.5cm]
    \item[Background formativo:] [formazione specifica pertinente al progetto]
    \item[Esperienza performativa:] [competenze esecutive e artistiche rilevanti]
    \item[Competenze tecniche:] [abilità specifiche in ambito tecnologico]
    \item[Capacità di ricerca:] [esperienze pregresse di ricerca documentate]
\end{description}

\subsection{Impatto e disseminazione previsti}
[Descrivere le ricadute attese nel campo degli studi musicali, le applicazioni pratiche e il potenziale di diffusione dei risultati]

\section{Distribuzione del lavoro nei tre anni}

\subsection{Primo anno}
\begin{itemize}[leftmargin=1cm]
    \item Ricerca bibliografica e archivistica approfondita
    \item Acquisizione di competenze tecniche specifiche
    \item Sviluppo del framework teorico di riferimento
    \item Prime sperimentazioni performative e metodologiche
\end{itemize}

\subsection{Secondo anno}
\begin{itemize}[leftmargin=1cm]
    \item Sviluppo operativo del progetto di ricerca
    \item Attivazione di collaborazioni con istituzioni partner
    \item Creazione di contenuti originali (composizioni, performance, software)
    \item Raccolta sistematica e analisi dei dati
\end{itemize}

\subsection{Terzo anno}
\begin{itemize}[leftmargin=1cm]
    \item Finalizzazione della ricerca e completamento degli obiettivi
    \item Realizzazione della produzione artistica conclusiva
    \item Stesura e revisione della dissertazione finale
    \item Disseminazione dei risultati attraverso pubblicazioni e performance
\end{itemize}

\section{Valorizzazione dell'heritage veneziano}
[Specificare in che modo il progetto valorizza il patrimonio musicale veneziano, si connette alla tradizione locale, contribuisce all'offerta culturale della città e integra prospettive interdisciplinari coerenti con l'identità del territorio]

\newpage

\section*{Bibliografia iniziale}
\begin{enumerate}
    \item Riferimento bibliografico completo 1
    \item Riferimento bibliografico completo 2
    \item Riferimento bibliografico completo 3
    \item Riferimento bibliografico completo 4
    \item Riferimento bibliografico completo 5
    \item ...
\end{enumerate}

\vspace{1cm}

\noindent\rule{\textwidth}{0.4pt}

\small
\textbf{Note:} Il presente progetto è stato strutturato in conformità ai requisiti del Bando per l'ammissione al Corso di Dottorato di ricerca AFAM in Musica, Performance e Innovazione Tecnologica, XLI ciclo, del Conservatorio di Musica "Benedetto Marcello" di Venezia. Il corpo del testo (esclusi frontespizio, bibliografia e apparati) rispetta il limite di 10.000 caratteri spazi compresi.

\end{document}
